\documentclass[10pt,a4paper]{amsart}
\usepackage[margin=19mm]{geometry}
\usepackage{amsmath,amssymb,mathtools}
\usepackage[scr=rsfs,cal=euler]{mathalfa}
\usepackage{xfrac}
\usepackage[unicode,hidelinks,bookmarks=false]{hyperref}
\newcommand{\ie}{i.e.\ }
\newcommand{\cf}{cf.\ }
\newcommand{\resp}{respectively\ }
\newcommand{\Subsection}{\subsection}
\providecommand{\nobreakdash}{\nobreak\hbox{-}}
\setlength{\emergencystretch}{2em}
\multlinegap0pt
\usepackage{bm}
\usepackage{bbm}
\usepackage{dsfont}
\usepackage{xspace}
\usepackage{cancel}
\usepackage{mathtools}
\usepackage{amsthm}
\makeatletter
\@ifundefined{theorem}{\newtheorem{theorem}{Theorem}}{}
\makeatother
\newcommand{\bmf}[1]{\mathbf{#1}}
\newcommand{\lbd}{\lambda}
\title[z-Classes in Finite Coxeter Groups]{z-Classes in Finite Coxeter Groups}
\author{Dilpreet Kaur, Uday Bhaskar Sharma}
\date{Source: arXiv:2401.05178v2 (CC BY 4.0)}
\begin{document}
\maketitle

\noindent\emph{Source and licence.} z-Classes in Finite Coxeter Groups. Version arXiv:2401.05178v2 (CC BY 4.0), distributed under the Creative Commons Attribution 4.0 International licence, as recorded in the arXiv licence metadata for this exact version and shown on the abstract page https://arxiv.org/abs/2401.05178v2. Full rights notice: licenses/arxiv-2401.05178-CC-BY-4.0.txt.\par\smallskip

\noindent\textit{Editorial scope.} Counted unit: the Theorem whose source label is \texttt{Th\_split\_CC\_in\_D\_n} in the article (source0.tex lines 598--600, source label \texttt{Th\_split\_CC\_in\_D\_n}), delivered together with the article's own complete original proof of it (source0.tex lines 602--662, one \texttt{proof} environment of 61 lines). No mathematics is written, repaired or re-derived: statement and proof are the article's own text line for line. Required context retained uncounted: none. Every definition, display and label of those lines is retained unchanged. Omitted from this package and not redistributed: 1--597, 601--601, 663--836. Nothing inside a retained excerpt is omitted or altered: every retained body is delivered exactly as the article's lines are written, so no omission receipt of any kind accompanies this package. Citations: the retained excerpts carry no citation command at all, so no bibliography entry of the article is transcribed and no reference is redistributed. Reference adaptation: none was needed and none was made; every reference inside the delivered unit resolves inside it, so no unresolved reference or citation is produced. Printed slips kept literal: no printed word, symbol, hypothesis or number of the counted statement or of its proof was altered. Layout and class: the article's own class is \texttt{amsart} with packages including xy, amssymb, latexsym, amsthm, amsmath, mathtools, dsfont, mathrsfs, bm, enumitem, color, makecell, biblatex, float; the wrapper is the lane's pinned \texttt{amsart} editorial wrapper at 10pt on A4, which restates the theorem-like environments the retained text needs and adds the article's own macro definitions that the retained text needs (\texttt{\textbackslash bmf}, \texttt{\textbackslash lbd}), with extra wrapper packages (bm, bbm, dsfont, xspace, cancel, mathtools). The delivered numbering is the unit's own. Assets: no retained span contains a figure environment, a graphics-inclusion command, a PDF-page inclusion, a table or a TikZ picture; no image file is copied into this package, the package declares no asset dependency and no image file is redistributed with it. Licence: version arXiv:2401.05178v2 of this article is distributed under the Creative Commons Attribution 4.0 International licence, as recorded in the arXiv licence metadata for this exact version and shown on the abstract page https://arxiv.org/abs/2401.05178v2. A full rights notice is delivered at licenses/arxiv-2401.05178-CC-BY-4.0.txt. Characters: every retained excerpt is pure ASCII, so the delivered engine, XeLaTeX with its default Latin Modern text font, reports no missing character.
\begin{theorem}\label{Th_split_CC_in_D_n}
Two split conjugacy classes in $D_n$ corresponding to the signed partition $\bmf{\bar{\lambda}} = m_1^{x_1}\cdots m_v^{x_v}$, where $2\leq m_1 < \cdots m_v\leq n$ are even, are $z$-conjugate if and only if there is a part $m_j = 4q_j + 2$ for some integer $q_j$, whose multiplicity $x_j$ is odd.
\end{theorem}

\begin{proof}
 Let $ \bmf{\bar{\lambda}}=m_1^{x_1}\cdots m_v^{x_v}$ with all even parts. Suppose there is $m_j = 4q_j + 2$, with multiplicity $x_j$ being odd. We shall reorder the parts of $\lbd$ so that $m_1 = 4q_1 + 2$, with odd multiplicity $x_1$. 

Let $\sigma = (1, \dots, m_1)\dots((x_1-1)m_1+1,\dots,x_1m_1)\dots(n-m_v+1,\dots,n)$
 and ${\bf a} = [{\bf 1}^n; \sigma]$ be an element of conjugacy class of $C_2 \wr S_n$ corresponding to signed partition $\bmf{\bar{\lambda}}.$
 Now consider the element ${\bf b} = [\bmf{-1}_{x_1m_1},\bmf{1}_{n-x_1m_1}; \sigma]$. We claim that ${\bf a}$ and ${\bf b}$ are not conjugate in $D_n.$ 

 
 Let ${\bf h} = [h_1,\ldots, h_n; \tau] \in C_2\wr S_n$ such that ${\bf hah^{-1}} = {\bf b}$, then $\tau \in Z_{S_n}(\sigma)$ and we have: $h_i h_{\sigma^{-1}(i)} = -1$, for $1 \leq i \leq x_1m_1$ and $h_ih_{\sigma^{-1}(i)} = 1$ for the remaining $i$'s. As $x_1$ and $\frac{m_1}{2}$ are odd, by a routine check, we can see that ${\bf h} \notin D_n$.

 As we are dealing with a split classes, we know that ${Z}_{D_n}({\bf a}) = {Z}_{C_2\wr S_n} ({\bf a})$. 

 Now, we know that 
 \begin{eqnarray*}
{Z}_{D_n}({\bf b}) &=& {Z}_{C_2 \wr S_n}({\bf b}) \\
&=& {Z}_{C_2\wr S_{r_1m_1}([{\bf -1}_{x_1m_1}; \sigma_{m_1.x_1}]}) \times \prod_{i = 2}^v {Z}_{C_2 \wr S_{x_im_i}}
([{\bf 1}_{x_im_i}; \sigma_{m_i.x_i}])\\
&=& {Z}_{C_2\wr S_{r_1m_1}([{\bf 1}_{x_1m_1}; \sigma_{m_1.x_1}]}) \times \prod_{i = 2}^v {Z}_{C_2 \wr S_{x_im_i}}
([{\bf 1}_{x_im_i}; \sigma_{m_i.x_i}])\\
&=& {Z}_{C_2 \wr S_n}({\bf a}) \\
&=& {Z}_{D_n}({\bf a})
\end{eqnarray*}
Thus ${\bf a}$ and ${\bf b}$ are $z$-conjugate in $D_n$.

Now, we consider the signed partition $\bmf{\bar{\lambda}} =m_1^{x_1}\cdots m_v^{x_v}$ corresponding to two split conjugacy classes in $D_n$, where all the parts of the type $m_j = 4q_j + 2$ have even multiplicity $x_j$. Here, we take ${\bf a} = [{\bf 1}_n; \sigma]$ and ${\bf b} = [{\bf -1}_2, {\bf 1}_{n-2}; \sigma]$. One can check that these two elements are not conjugate in $D_n.$

We know that: 
\begin{eqnarray*}
 {Z}_{D_n}({\bf b}) &=& {Z}_{C_2 \wr S_n}({\bf b}) \\
&=&{Z}_{C_2\wr S_{x_1m_1}}([{\bf -1}_2,{\bf 1}_{x_1m_1-2}; \sigma_{m1.x1}]) \times \prod_{i = 2}^v {Z}_{C_2 \wr S_{x_im_i}}([{\bf 1}_{x_im_i}; \sigma_{m_i.x_i}])
\end{eqnarray*}

Suppose ${\bf a}$ and ${\bf b}$ are $z$-conjugate in $D_n.$ We have ${\bf g}{Z}_{D_n}({\bf a}){\bf g}^{-1} = {Z}_{D_n}({\bf b})$ for some ${\bf g} \in D_n$
This implies
\begin{eqnarray*}{\bf g}Z({Z}_{D_n}({\bf a})){\bf g}^{-1}& =& Z({Z}_{D_n}({\bf b}))\\
\Rightarrow{\bf g}\prod_{i=1}^v\langle[{\bf 1}_{x_im_i}; \sigma_{m_i.x_i}],[{\bf -1}^n; (1)]\rangle{\bf g}^{-1} &=& \langle[{\bf -1}_2,{\bf 1}_{x_1m_1-2}; \sigma_{m_1.x_1}],[{\bf -1}_{x_1m_1}, (1)]\rangle\\&~& \times \prod_{i=2}^v\langle[{\bf 1}_{x_im_i}; \sigma_{m_i.x_i}],[{\bf -1}_{n-m_1x_1}; (1)]\rangle
\end{eqnarray*}
Therefore, we get ${\bf g}[{\bf 1}_n; \sigma]{\bf g}^{-1} = [{\bf -1}_2, {\bf 1}_{n-2}; \sigma]^w={\bf b}^w,$ where $(w, o(\bmf{a})) = 1$, in particular, $(w, m_1) = 1$. 

Now, let ${\bf g}=[g_1,\ldots,g_n; \tau]$, where $\tau\sigma\tau^{-1} = \sigma^w,$ 
and we compute 

${\bf b}^w = [-1,{\bf 1}_{w-1},-1,{\bf 1}_{n-w-1}; \sigma^w].$

On equating ${\bf gag}^{-1} = {\bf b}^w$, for $1 \leq i \leq m_1,$ we get
\begin{eqnarray*}
g_1g_{m_1 -w +1} &=& -1\\
g_2g_{m_1 -w + 2} &=& 1\\
\vdots &=& \vdots\\
g_{w+1}g_1 &=& -1\\
\vdots &=&\vdots\\
g_{m_1}g_{m_1 -w} &=& 1
\end{eqnarray*}

We recall that $\sigma_{m_1}^w$ denotes the $(12\cdots m_1)^w.$ As $\sigma_{m_1}^w= (1\sigma_{m_1}^w(1) \cdots \sigma_{m_1}^{w(m_1-1)}(1)),$ we have
$g_{\sigma_{m_1}^{w(m_1-1)}(1)} =g_{\sigma_{m_1}^{w(m_1-2)}(1)} = \cdots = g_{\sigma_{m_1}^w(1)} = -g_1$. 

But  $\{\sigma_{m_1}^w(1), \ldots, \sigma_{m_1}^{w(m_1-1)}(1)\} = \{2,3,\ldots, m_1\}$ we get that $-g_1 = g_2 = \cdots = g_{m_1}.$

For $m_1 + 1 \leq j \leq n$, we have $g_{j}g_{\sigma^{-w}(j)} = 1$. Thus, the element ${\bf g}$ is of the form $${\bf g}= [-d^{(1)}_{1},{\bf d^{(1)}_{1}}_{m_1 - 1},{\bf d^{(1)}_2}_{m_1},\ldots,{\bf d^{(1)}_{x_1}}_{m_1},\cdots, {\bf d^{(v)}_{x_v}}_{m_v}; \tau] ,$$ where $d^{(i)}_j,= \pm 1$, for $1 \leq i \leq v$ and $1 \leq j \leq x_i$. The total number of $-1$ in ${\bf g}$ is either $2q + 1$ or $2q + (m_1 -1)$, which is odd. This implies ${\bf g} \notin D_n$. Hence ${\bf a}$ and ${\bf b}$ are  not $z$-conjugate in $D_n.$ 
\end{proof}

\end{document}
